% ===================================================================
% Major-revision supplement: composite-objective Stage-1 (v3) results.
% Written to be \input from main_is.tex. Honest framing: the gain comes
% from the COMPOSITE optimizer-facing training objective, NOT from
% differentiable projection in the loop (post-hoc projection ties it).
% ===================================================================

\section{Composite-Objective Prior Training and Ablations}
\label{sec:mr-composite}

The preceding sections fixed the projection operator and treated the Stage-1
prior as secondary. We now ask the converse, reviewer-facing question directly:
\emph{when the prior is trained on an optimizer-facing objective rather than on
histogram reconstruction, where does the gain come from, and is it the
calibration step or the objective?} We retrain the prior generator (a small
permutation-invariant set encoder over the stale-quantile and feedback tokens) on
a composite loss evaluated on held-out future predicates,
\[
  \mathcal{L} = w_{\text{fut}}\,\mathrm{QErr}_{\log}\!\big(\Pi(p_\theta),\sigma_{\text{fut}}\big)
              + w_{\text{join}}\,\mathrm{Regret}_{\text{FJ}}\!\big(\Pi(p_\theta)\big)
              + w_{\text{cons}}\,\rho_{\text{fb}}(p_\theta)
              + w_{\text{reg}}\,\lVert p_\theta - p_{\text{fresh}}\rVert_1,
\]
where $\Pi$ is the feedback projection, $\mathrm{QErr}_{\log}$ the held-out
future-predicate log-Q-error, $\mathrm{Regret}_{\text{FJ}}$ a differentiable
FactorJoin bin-mass join regret, $\rho_{\text{fb}}$ the in-window feedback
residual, and the last term a light reconstruction regularizer toward the
post-drift marginal. Training uses a sparse-to-dense feedback curriculum and a
domain-randomized drift mixture; all comparisons below hold the data, capacity,
and projection operator fixed across variants.

\subsection{Ablation: the gain is the objective, not in-the-loop calibration}
\label{sec:mr-ablation}

Table~\ref{tab:mr_ablation} decomposes the objective. Each row is gate-evaluated
identically: the model's prior is calibrated by the \emph{real} projection on
held-out future predicates, and reported as the percentage improvement over the
strongest classical prior (STHoles) at three feedback budgets.

\begin{table}[t]
  \centering\small
  \caption{Stage-1 objective ablation (held-out future-predicate Q-error after the
  same real projection; \% improvement over STHoles-init, higher is better). All
  variants share data, capacity, and projection. \emph{Reconstruction-only} is the
  floor; the composite objective drives the gain; \emph{raw-train + post-hoc
  projection} matches \emph{full in-loop}, so differentiable projection
  \emph{in the loop} is not the source of the gain. The dense-feedback ($K{=}16$)
  column shows the reconstruction/consistency regularizers, not in-loop
  projection, are what prevent collapse.}
  \label{tab:mr_ablation}
  \setlength{\tabcolsep}{6pt}
  \begin{tabular}{l rrr}
    \toprule
    Stage-1 objective variant & $K{=}2$ & $K{=}6$ & $K{=}16$ \\
    \midrule
    Reconstruction-only (old objective) & $-1.2$ & $-1.9$ & $-7.8$ \\
    Future-loss only (in-loop proj.)    & $+8.8$ & $+7.0$ & $+0.5$ \\
    \;\;+ FactorJoin regret             & $+8.0$ & $+5.9$ & $-0.0$ \\
    \;\;+ feedback residual             & $+9.8$ & $+7.4$ & $+3.4$ \\
    Raw-train + \emph{post-hoc} projection & $+9.4$ & $+9.0$ & $+7.8$ \\
    \textbf{Full (in-loop projection)}  & $+8.7$ & $+7.9$ & $+7.8$ \\
    \bottomrule
  \end{tabular}
\end{table}

Three conclusions follow, and we state the one that cuts against a tempting
narrative explicitly. First, \emph{reconstruction-only fails at every budget}
(it is worse than the classical priors), so an optimizer-facing objective is
necessary---training a prior to look like the post-drift histogram does not make
it a better optimizer input. Second, the composite objective is what creates the
gain, and the reconstruction and consistency regularizers are essential at dense
feedback: future-loss alone, and future-loss${+}$join, collapse to a tie at
$K{=}16$, and only adding the regularizers restores the dense-feedback margin.
Third---and contrary to a ``calibration-in-the-loop'' reading---\emph{training
the prior through the projection operator is not the source of the gain}:
``raw-train ${+}$ post-hoc projection'' (the loss never sees the projection;
projection is applied only at evaluation) matches ``full in-loop'' to within
$0.5\%$ at every budget, and this holds across two training seeds on the full
drift range (post-hoc $+6.7/+7.6/+6.9$ and $+6.8/+6.7/+7.1$; in-loop
$+6.8/+7.0/+6.2$ and $+6.8/+6.0/+4.3$ at $K{=}2/6/16$). We therefore do not claim
that differentiable projection in the loop is necessary; it is a convenient but
inessential training device, and post-hoc projection---simpler to
deploy---performs identically. The contribution is the \emph{composite
optimizer-facing training objective}, not the in-loop calibration.

\subsection{Why sparse feedback helps and dense feedback ties: a rank argument}
\label{sec:mr-mechanism}

The regime structure is not an empirical accident; it follows from the rank of
the feedback-constraint system. An equi-depth marginal with $B$ buckets has
$B{-}1$ free interior quantile boundaries. Each feedback predicate constrains the
CDF at its boundary value(s); after projection, the marginal is pinned only along
the span of those constraints, and the orthogonal complement (the nullspace) is
left free for the prior to fill. Table~\ref{tab:mr_rank} reports, per budget $K$,
the mean independent-constraint rank and the residual free degrees of freedom
over the held-out cases ($B{=}10$, so $B{-}1{=}9$).

\begin{table}[t]
  \centering\small
  \caption{Feedback-constraint rank and residual free degrees of freedom (DOF) of
  the marginal vs feedback budget $K$ (marginal DOF $=B{-}1=9$). Free DOF is the
  nullspace dimension the projection leaves for the prior to complete; once
  feedback saturates the DOF the marginal is pinned and the prior washes out.}
  \label{tab:mr_rank}
  \setlength{\tabcolsep}{8pt}
  \begin{tabular}{r rr r}
    \toprule
    $K$ & mean rank & mean free DOF & marginals pinned \\
    \midrule
    2  & 2.32 & 6.68 & 0\% \\
    4  & 4.65 & 4.35 & 0\% \\
    6  & 6.95 & 2.05 & 6\% \\
    8  & 8.71 & 0.29 & 73\% \\
    12 & 8.98 & 0.02 & 98\% \\
    16 & 8.98 & 0.02 & 98\% \\
    \bottomrule
  \end{tabular}
\end{table}

The free DOF collapses from $6.68$ at $K{=}2$ to $\approx 0$ by $K{\geq}12$, with
$73\%$ of marginals already fully pinned at $K{=}8$. This maps directly onto the
measured behavior: the learned prior wins where free DOF is large (sparse
feedback), the margin shrinks as the nullspace closes ($K{=}8$), and at dense
feedback the marginal is pinned by the constraints so the projection determines
it regardless of the prior---OASIS converges to ISOMER. Sparse feedback is
precisely the regime where the optimizer-relevant completion the prior supplies
in the nullspace cannot be recovered from the feedback alone.

\subsection{The Router does not use oracle information}
\label{sec:mr-router}

The Router selects among candidate calibrated marginals using \emph{only} the
in-window feedback residual---the mean absolute discrepancy between each
candidate's estimate and the observed selectivity on the current feedback
window. It never observes test predicates, true cardinalities, runtimes, or plan
shapes; the selection signal is exactly what a deployed system already has. Its
per-budget choices follow the rank mechanism above: it selects the calibrated
learned prior at $K\in\{2,4,6,8\}$ (where free DOF is available) and falls back
to ISOMER at $K{=}16$ (where the marginal is pinned and the candidates coincide).
The selection is plan-shape-safe: at dense planner feedback the routed output
introduces new plan deviations in $1.1\%$ of queries, identical to ISOMER and
below the raw prior's $2.3\%$.

\subsection{Full-stack TPC-H safety: reproduces fresh-statistics behavior}
\label{sec:mr-tpch}

We re-run the TPC-H SF10 DML-drift study with the composite-objective prior over
three refresh-stream seeds, and report not only scan accuracy but the runtime
distribution and plan-safety statistics a deployment reviewer cares about
(Table~\ref{tab:tpch_runtime}). Because warm-cache runtime is determined by the
plan the injected statistics induce, we time each \emph{distinct} (query, plan)
once at full repetition and assign it to every method sharing that plan; this
removes redundant timing and surfaces directly that calibrated and fresh
statistics induce the same plans on these queries.

The accuracy story is stable across seeds. Stale statistics are badly wrong on the
drifted date columns (scan Q-error $5.23\pm0.24$), and calibration repairs them to
fresh quality (OASIS $2.02$ and the Router $2.01$, matching fresh $2.01$). The
runtime story is the deliberately bounded one. Calibrated statistics reproduce the
\emph{fresh}-statistics runtime to within a few percent (geometric-mean
time-to-fresh ratio $0.96$ for OASIS and $1.01$ for the Router), because they
induce the fresh plans. The \emph{stale} wall-clock baseline, however, is not a
reliable target: on this configuration the stale mis-estimate happens to pick a
cheaper-to-run plan, so calibrated (and fresh) statistics run $\approx 1.5\times$
\emph{slower} than stale (time-to-stale $1.49$ for OASIS, $1.57$ for the Router,
$1.56$ for fresh itself), and a handful of queries per seed exceed the $15\%$
regression band \emph{relative to stale}---but they regress relative to
\emph{fresh} too, i.e.\ they are the price of the correct plan, not a calibration
artifact. In a larger-drift configuration the same stale mis-estimate instead
produced a plan that ran $1.56\times$ \emph{slower} than fresh, so the direction of
the stale-vs-fresh gap is configuration-dependent. The deployment-relevant
invariant is that calibrated statistics track fresh statistics in both accuracy and
runtime, whereas stale statistics make the wall-clock outcome unpredictable in
either direction. Maintaining these statistics requires one forward pass plus the
projection ($\approx 1.13$\,ms per column, Section~\ref{sec:overhead}) and avoids
the full-table \texttt{ANALYZE} scan that an SF10 refresh would otherwise incur.
The claim remains the bounded one: calibrated statistics \emph{reproduce
fresh-statistics behavior} without the scan; we make no broad
runtime-superiority claim over leaving statistics stale.

\subsection{Baseline fairness}
\label{sec:mr-fairness}

All methods are compared under identical conditions. Every method consumes the
\emph{same} feedback window (the same $K$ predicate observations) and starts from
the \emph{same} stale statistics; ISOMER, STHoles, QuickSel-H, and the learned
prior are all projected onto the \emph{same} feedback constraints by the same
hard projection, and all are injected into PostgreSQL through the \emph{same}
\texttt{pg\_statistic} path and evaluated on the \emph{same} held-out predicate
split. No method receives oracle access: the Router and all gates see only
in-window feedback residuals, never test predicates, true cardinalities, or
runtimes; ``fresh'' is computed from the post-drift data purely as an accuracy
upper bound and is not available to any deployed method. The learned prior and
the classical priors are evaluated at every feedback budget under the identical
protocol, so the regime comparison is like-for-like.

\subsection{Scope and what we do not claim}
\label{sec:mr-scope}

We are deliberate about the boundary of these results. (i) OASIS repairs
\emph{single-column} marginal statistics; correlated multi-column predicates are
not solved here, and where joint-cardinality error is dominated by cross-column
dependence an external dependence model remains necessary. (ii) The FactorJoin
kernel is a \emph{join-aware proxy} that exercises how a corrected single-column
key histogram propagates through a bilinear join estimate; it is not a full
optimizer replacement. (iii) At dense feedback the marginal is pinned by the
constraints (Section~\ref{sec:mr-mechanism}), and we explicitly do \emph{not}
claim the learned prior beats projection there---it provably cannot, and ties
ISOMER by design. (iv) Fresh statistics remain the accuracy upper target the
method reproduces, not an object it replaces. The contribution is the
composite optimizer-facing training objective for the prior, together with the
regime-aware calibration and routing that make imperfect marginals safe to
consume; in-loop differentiable projection is an inessential training device
(Section~\ref{sec:mr-ablation}).

\subsection{Statistical significance and robustness}
\label{sec:mr-significance}

The improvements are reported over multiple seeds and held-out splits rather than
single runs: the PostgreSQL planner study aggregates twelve configurations
(two drift directions $\times$ two table sizes $\times$ three seeds), the TPC-H
study three refresh-stream seeds, and the objective ablation two training seeds
on the full eight-intensity drift range (the in-loop vs.\ post-hoc equivalence
holds on both). Robustness is probed along every axis a reviewer would test:
feedback budget $K$ (the rank sweep, Section~\ref{sec:mr-mechanism}), drift
severity (the per-intensity $q$ curves), drift \emph{family} via
leave-one-family-out generalization (the prior is trained on five
deployment-inspired drift families and evaluated on a held-out sixth, where it
still beats the classical priors and ties ISOMER), and feedback noise (graceful
degradation to $10\%$ multiplicative noise). The failure modes are reported, not
hidden: the raw prior is actively harmful in the FactorJoin kernel and on the
append-only NASA trace, which is exactly why the projection and residual gating
are part of the method; and at extreme sparsity on a never-seen drift family the
learned prior reaches parity with ISOMER rather than strictly dominating it.
